\chapter{Conclusioni}\label{conclusioni}

\section{Raggiungimento degli
obiettivi}\label{raggiungimento-degli-obiettivi}

Il presente lavoro di tesi si è posto l'obiettivo di
superare la barriera tecnica tra il linguaggio naturale umano e la
rigidità sintattica delle basi di dati relazionali. L'attività di
ricerca e sviluppo si è concretizzata nella realizzazione di un
framework basato su Python, strutturato per agire come un intermediario
intelligente capace di orchestrare l'intero ciclo di vita di una
richiesta text-to-SQL.

L'obiettivo principale è stato perseguito attraverso l'implementazione
di una serie di moduli interconnessi, la cui efficacia è stata
verificata in fase sperimentale.

Di seguito vengono analizzate
nel dettaglio le soluzioni tecniche ideate per il raggiungimento
degli obiettivi.

L'obiettivo di estrarre e formalizzare la conoscenza del database è
stato raggiunto attraverso lo sviluppo della classe
\texttt{SchemaOrchestrator}. Questo componente gestisce l'acquisizione
del metadato tecnico operando secondo due direttrici: l'estrazione
diretta tramite connessione al database (\texttt{db\_conn}) o
l'elaborazione di descrizioni testuali fornite dall'utente tramite Large
Language Models. Il risultato di tale processo è la creazione di un
oggetto \texttt{Schema}, che funge da rappresentazione strutturata e
centralizzata del patrimonio informativo del database, fornendo una
base solida per le successive fasi di generazione.

Per ovviare ai limiti della finestra di contesto dei modelli linguistici, è
stato implementato un sistema di gestione del contesto basato su
architettura RAG. Il componente \texttt{SchemaStore} agisce come un
database vettoriale in cui lo schema viene frammentato e immagazzinato
sotto forma di documenti. Grazie all'integrazione di un
retriever, il sistema è in grado di analizzare la richiesta in
linguaggio naturale dell'utente e isolare dinamicamente solo le porzioni
di schema (tabelle e relazioni) realmente pertinenti al contesto della
query. Questo approccio non solo ottimizza il consumo di token, ma
è in grado anche di ridurre il ``rumore'' informativo, migliorando
l'accuratezza del modello.

La flessibilità del sistema nella generazione del codice SQL è fornita
dal \texttt{PromptBuilder}. Questo modulo è responsabile della
costruzione di prompt dinamici che si adattano a molteplici variabili:
la tipologia di sorgente dati e le specifiche necessità sintattiche del
dialetto SQL di destinazione. Il \texttt{PromptBuilder} integra in un
unico template la richiesta dell'utente, il contesto estratto dallo
\texttt{SchemaStore} e, in fase di raffinamento, i feedback derivanti
dai fallimenti precedenti, consentendo una comunicazione coerente e
 contestualizzata con il modello linguistico.

Il requisito di robustezza della query prodotta è stato soddisfatto
mediante un ciclo iterativo di validazione e raffinamento gestito dal
\texttt{QueryOrchestrator}. Ogni query generata viene sottoposta a una
duplice verifica: una validazione sintattica immediata (tramite librerie
di parsing come \texttt{sqlglot}) e, ove possibile, una validazione
semantica attraverso l'esecuzione reale sul database. In caso di errore
(sia esso di sintassi o di runtime), interviene la logica di
\texttt{LLMFeedback}, che analizza l'eccezione e produce un'indicazione
correttiva che viene reiniettata nel prompt per il tentativo successivo.
Questo meccanismo di ``auto-correzione'' permette al sistema di imparare
dai propri errori all'interno della medesima sessione, incrementando
significativamente il tasso di successo finale delle interrogazioni
generate.

L'ultimo requisito essenziale, finalizzato a rendere il framework
fruibile, modulare e facilmente integrabile all'interno di ecosistemi
applicativi di terze parti, è stato conseguito attraverso la
progettazione e l'implementazione di uno strato di servizi web basato
sul framework \textbf{FastAPI}. Sfruttando le caratteristiche native di
asincronia, validazione dei tipi e generazione automatica della
documentazione offerte da questa tecnologia, l'intera logica di business
è stata esposta tramite un'architettura ad \textbf{API RESTful} solida e
ad alte prestazioni.

La modularità del sistema si riflette a livello infrastrutturale
nell'adozione di una strategia di routing distribuita. Il punto di
ingresso dell'applicazione (\texttt{main.py}) garantisce la massima
interoperabilità frontend e orchestra l'inclusione di router dedicati ed
isolati, nello specifico \texttt{schema} e \texttt{query}. Questo
approccio assicura una netta separazione delle responsabilità e una
manutenibilità ottimale del codice sorgente. \textbf{Il router
\texttt{schema.py}} centralizza gli endpoint dedicati alla gestione del
ciclo di vita dei metadati. Esso espone le rotte per l'estrazione
diretta delle strutture relazionali da istanze MySQL remote
(\texttt{/schema/mysql}), per la generazione sintetica di schemi a
partire da testo non strutturato mediante modelli linguistici
(\texttt{/schema/text}) e per il raffinamento iterativo delle
definizioni esistenti (\texttt{/schema}). \textbf{Il router
\texttt{query.py}} si occupa invece di governare i flussi operativi di
traduzione del linguaggio naturale. Al fine di massimizzare la
flessibilità del framework in contesti di test o di produzione offline,
sono stati ingegnerizzati due flussi logici distinti: l'endpoint
\texttt{/queries/text}, deputato alla sola generazione pura del codice
SQL isolando la componente LLM da interazioni esterne, e l'endpoint
\texttt{/queries/mysql}, il quale estende il precedente paradigma
abilitando l'interconnessione con il DBMS, l'esecuzione della query
transazionale e la serializzazione formale dei record risultanti.

In sintesi, la traduzione dei requisiti iniziali in un'architettura
software funzionante ha dimostrato che l'integrazione tra modelli di
linguaggio avanzati e sistemi di recupero delle informazioni rappresenta
una soluzione robusta per democratizzare l'accesso ai dati, riuscendo ad ottenere
al contempo sia la sicurezza che la correttezza formale delle interrogazioni
eseguite.

\section{Vantaggi e svantaggi}\label{vantaggi-e-svantaggi}

La valutazione sistematica delle prestazioni ha permesso di delineare
con precisione i benefici introdotti dall'architettura proposta, nonché
i vincoli tecnologici che ancora ne limitano l'efficacia in scenari di
produzione complessi.

Il principale valore aggiunto del sistema risiede nella sua capacità di
auto-correzione iterativa, poiché l'integrazione della modalità
\emph{db\_conn} trasforma il framework da un semplice traduttore statico
a un agente dinamico in grado di avvicinarsi sempre di più all'obiettivo finale.

Questa architettura favorisce una netta riduzione delle allucinazioni
sintattiche, come dimostrato da una padronanza grammaticale del
linguaggio SQL pressoché assoluta; l'assenza di errori di sintassi
suggerisce infatti la validità dei template di prompt e dei vincoli
strutturali imposti ai modelli portando la media generale dei successi da \(23,78\%\) a \(45,33\%\). Al contempo, si rileva un'elevata
efficacia del feedback loop: per modelli specializzati come
SQLCoder il sistema agisce come un catalizzatore di precisione, poiché
anche laddove la prima generazione fallisce, il meccanismo di cattura
degli errori di runtime e la loro successiva re-iniezione nel prompt
permettono di ``raddrizzare'' la query, garantendo il successo finale (da \(16,39\%\) a \(25,39\%\)).

L'analisi ha inoltre evidenziato la versatilità dei modelli open-source,
mostrando come soluzioni quali Codestral e Qwen3-Coder-Next siano in grado
di competere con architetture proprietarie del calibro di GPT-4o in
contesti zero-shot, offrendo un'alternativa valida, economica e
orientata alla privacy per l'interrogazione di schemi non eccessivamente
densi. Per i primi due si osserva \(48,46\%\) e \(55,77\%\) rispettivamente 
mentre per i due successivi \(,54\%\) e \(61,54\%\). Di conseguenza, il 
sistema raggiunge l'obiettivo della democratizzazione del dato, abbattendo la barriera tecnica per
database con complessità moderata e permettendo l'estrazione di
informazioni complesse, quali join e aggregazioni, senza la necessità di
conoscere lo specifico dialetto SQL.

Nonostante i successi prestazionali, emergono criticità strutturali
legate principalmente alla gestione del contesto e alla latenza
operativa.

Il limite strutturale primario è rappresentato dalla sensibilità alla
densità dello schema, specialmente nei database ampio come
\emph{california\_schools} con 89 colonne totali, dove un numero elevato di attributi per
singola tabella tende a saturare il contesto del modello. In questi
scenari, il sistema fatica a distinguere tra termini semanticamente
simili, producendo di conseguenza errori di runtime legati a nomi di
colonne errati.

A questa barriera architetturale si aggiunge un'insidiosa vulnerabilità
legata al fenomeno del Value Mapping, emersa dall'ispezione
approfondita dei log di sistema su domini caratterizzati da un limitato
punteggio sulla complessità generale ma da un severo tasso di fallimento
logico (quali \emph{real\_estate\_properties}, \emph{card\_games},
\emph{toxicology}, \emph{thrombosis\_prediction} e \emph{superhero}). In
questi contesti, la corretta esecuzione della query richiede
l'applicazione di filtri restrittivi nella clausola \texttt{WHERE} su
attributi di tipo \texttt{VARCHAR}. Tuttavia, i valori categoriali o
testuali esatti necessari a soddisfare il predicato algebrico non sono
specificati nella richiesta in linguaggio naturale, né sono estraibili
dai metadati strutturali. Privi di un meccanismo di indicizzazione dei
contenuti del database, i modelli eseguono ripetuti tentativi alla
cieca inserendo costanti letterali plausibili ma fattualmente errate;
in questo modo il DBMS esegue la query con successo azzerando i runtime
error, ma restituisce un set di dati vuoto. Tale dinamica provoca una
massiccia traslazione dei fallimenti verso la dimensione semantica,
esaurendo inutilmente il budget di raffinamento iterativo.

A questo si sommano i limiti su latenza e costi computazionali, poiché
l'adozione di un approccio multi-tentativo, sebbene efficace per
incrementare la precisione, penalizza in modo significativo la velocità
di risposta; modelli come GPT-5-mini, pur risultando accurati,
presentano tempi di inferenza che possono superare i 50 secondi su
database complessi, rendendo l'esperienza utente meno fluida in contesti
operativi in tempo reale.

Si osserva inoltre una fragilità semantica in assenza di connessione
nella modalità \emph{text}, in cui il sistema dipende totalmente dalla
qualità della documentazione dello schema. Senza l'ausilio della
connessione diretta al database, la capacità di ricostruire relazioni
logiche profonde cala, come dimostrato dal caso limite del
database \emph{orchestra}, evidenziando una dipendenza critica
dall'introspezione dinamica, possedendo un delta tra le due modalità di \(43,33\%\). 
Infine, l'analisi ha rilevato un'efficacia
limitata del processo di raffinamento sui modelli generalisti;
alcune architetture, pur essendo avanzate, non traggono benefici
significativi dai tentativi aggiuntivi, mostrando una correlazione tra
tentativi e successo vicina allo zero (esempio: Qwen3-Coder-Next). Questo fenomeno suggerisce che
per tali agenti l'errore iniziale sia spesso di natura logica e,
pertanto, non facilmente correggibile attraverso il solo feedback
testuale fornito dal database.

\section{Prospettive future}\label{prospettive-future}

Sulla base delle criticità emerse in fase di valutazione sperimentale,
sono state identificate diverse direttrici di ricerca e ottimizzazione
ingegneristica volte a estendere le capacità e la robustezza del
framework proposto.

Per ovviare al crollo delle prestazioni riscontrato nei database ad alta
densità (caratterizzati da un elevato numero di attributi per singola
tabella), una naturale evoluzione del sistema prevede l'introduzione di
un modulo di Reranking integrato nello \texttt{SchemaStore}, ad
esempio mediante l'impiego di modelli Cross-Encoder. Questo approccio
consentirebbe un filtraggio del contesto a due stadi: un primo recupero
preliminare basato sulla similarità coseno degli embedding vettoriali,
seguito da una fase di riordinamento semantico fine. Tale strategia
garantirebbe l'iniezione nel prompt del solo sottoinsieme di colonne
strettamente necessarie alla risoluzione del quesito, riducendo il consumo di token nella finestra di contesto e
azzerando quasi totalmente gli errori di Schema Linking.

Le criticità emerse nell'ambito del value mapping hanno
evidenziato come, in assenza di informazioni puntuali sul contenuto dei
record, i modelli linguistici non dispongano di elementi sufficienti per
formulare correttamente i predicati di filtraggio (clausole
\texttt{WHERE}). Una promettente linea di sviluppo consiste
nell'ingegnerizzazione di un meccanismo di campionamento dinamico dei
dati. All'atto della generazione del prompt, il sistema potrebbe
interpellare direttamente il database target eseguendo query di tipo
\texttt{SELECT\ DISTINCT} sulle colonne di tipo testuale
(\texttt{VARCHAR} o \texttt{TEXT}) rilevanti. Valutando preventivamente
la cardinalità e i valori effettivi di tali attributi, l'LLM verrebbe
arricchito di una consapevolezza fattuale sul contenuto informativo del
database. Questo approccio eliminerebbe la necessità di effettuare
tentativi iterativi basati su valori errati o inesistenti, incrementando
in modo significativo il tasso di successo esecutivo della query finale.